\documentclass[10pt]{article}
\usepackage[T1]{fontenc}
\usepackage[utf8]{inputenc}
\DeclareUnicodeCharacter{2013}{--}
\usepackage{lmodern}
\usepackage{amsmath,amssymb,amsthm,mathtools}
\usepackage{geometry,multicol}
\usepackage{xurl}
\ifdefined\screenedition
  \geometry{paperwidth=1200pt,paperheight=700pt,margin=35pt,columnsep=18pt}
  \newenvironment{notesbody}{\begin{multicols}{2}\raggedcolumns}{\end{multicols}}
\else
  \geometry{letterpaper,margin=0.85in}
  \newenvironment{notesbody}{}{}
\fi
\usepackage[table]{xcolor}
\usepackage{graphicx}
\graphicspath{{../../LaTeX_and_pdfs/}}
\usepackage{tikz}
\usetikzlibrary{matrix,arrows,arrows.meta,cd,calc,decorations.pathmorphing,backgrounds}
\usepackage{listings}
\usepackage[version=3]{mhchem}
\usepackage[colorlinks=true,linkcolor=blue,urlcolor=blue,citecolor=blue]{hyperref}
\definecolor{darkgreen}{rgb}{0,0.4,0}
\lstset{language=Python,
 frame=bottomline,
 breaklines=true,
 basicstyle=\scriptsize,
 identifierstyle=\color{blue},
 keywordstyle=\bfseries,
 commentstyle=\color{darkgreen},
 stringstyle=\color{red},
 }
\theoremstyle{plain}
\newtheorem{theorem}{Theorem}
\newtheorem{lemma}{Lemma}
\newtheorem{proposition}{Proposition}
\newtheorem{corollary}{Corollary}
\newtheorem{axiom}{Axiom}
\theoremstyle{definition}
\newtheorem{definition}{Definition}
\newtheorem{remark}{Remark}
\newenvironment{claim}[1]{\par\noindent\underline{Claim:}\space#1}{}
\newenvironment{claimproof}[1]{\par\noindent\underline{Proof:}\space#1}{\hfill $\blacksquare$}
\newenvironment{correction}[1]{%
  \par\medskip\noindent\textbf{Correction (#1).}\space\ignorespaces
}{\par\medskip}
\newcommand{\questionhead}[1]{%
  \bigskip\bigskip\noindent{\small\bf Question #1.}\bigskip}
\newcommand{\problemhead}[1]{\noindent{\small\bf Problem #1.}}
\newcommand{\exercisehead}[1]{\smallskip\noindent{\small\bf Exercise #1.}}
\newcommand{\solutionhead}[1]{\noindent{\small\bf Solution #1.}}
\setlength{\emergencystretch}{3em}
\allowdisplaybreaks
\author{Ernest Yeung}
\date{12 October 2026}

\usepackage[makeroom]{cancel}
\usepackage{caption,pdfpages,pgfplots,pifont}
\pgfplotsset{compat=1.18}
\newcommand\dblwr{\wr\mkern-2mu\wr}
\newcommand{\seminorm}[1]{\mathopen{\dblwr}#1\mathclose{\dblwr}}
\theoremstyle{definition}
\newtheorem*{exercise}{Exercise}

\makeatletter
\renewcommand{\@pnumwidth}{2.5em}
\renewcommand{\@tocrmarg}{3.5em}
\makeatother
% Large tensor-network diagrams retain their original scale on wide pages.
\newcommand{\widetopic}[1]{%
  \end{notesbody}\clearpage
  \begingroup
  \newgeometry{left=35pt,top=35pt,textwidth=1130pt,textheight=630pt}
  \paperwidth=1200pt\paperheight=700pt
  \pdfpagewidth=1200pt\pdfpageheight=700pt
  \input{#1}\clearpage
  \endgroup
  \restoregeometry
  \begin{notesbody}
}
